\subsection{Dataset and signal representation}

The analysis used pretreatment EEG from 21 adults with treatment-resistant depression enrolled in a fully remote, multisite, double-blind, randomized sham-controlled home-based tDCS trial for major depressive disorder \citep{woodham2025home}. The analyzed cohort contained 7 remission cases and 14 non-remission cases in the patient-level outcome labels, with 18 female participants and a mean age of 37.1 $\pm$ 9.7 years. Table~\ref{tab:cohort-summary} summarizes the cohort, acquisition, and sweep configuration.

\begin{table}[H]
  \centering
  \caption{Cohort, acquisition, and sweep configuration summary. Trial descriptors derive from the parent home-based tDCS study \citep{woodham2025home}; acquisition totals and sweep settings derive from the analyzed EEG exports and current experiment configuration.}
  \label{tab:cohort-summary}
  \small
  \begin{tabular}{>{\raggedright\arraybackslash}p{0.42\linewidth}>{\raggedright\arraybackslash}p{0.50\linewidth}}
\toprule
Characteristic & Value \\
\midrule
\multicolumn{2}{l}{\textbf{Parent trial and cohort descriptors}} \\
\addlinespace[0.2em]
Parent trial & Fully remote, multisite, sham-controlled home-based tDCS trial for MDD \citep{woodham2025home} \\
Participants & 21 \\
Outcome groups & 7 remission, 14 non-remission \\
Sex (female/male) & 18/3 \\
Age (years, mean $\pm$ SD) & 37.1 $\pm$ 9.7 \\
\midrule
\multicolumn{2}{l}{\textbf{EEG acquisition and export descriptors}} \\
\addlinespace[0.2em]
Rest condition & Eyes-closed resting state \\
EEG montage & 4 channels (AF7, AF8, TP9, TP10) \\
Channel coverage & Frontal (AF7, AF8), temporal (TP9, TP10) \\
Sampling rate & 256 Hz \\
Recording structure & 10 min total; single file or two 5 min files \\
Available export segment length & 10 s (2,560 samples/channel) \\
Total exported windows & 1,203 \\
Remission windows & 393 (32.7\%) \\
Non-remission windows & 810 (67.3\%) \\
Class ratio & 2.06:1 (non-remission:remission) \\
\midrule
\multicolumn{2}{l}{\textbf{Current sweep methodology settings}} \\
\addlinespace[0.2em]
Sweep window sizes & 2, 4, 6, 8, 10 s \\
Inner-k sweep & 1, 5, 10, 15, 20, 25, 30, 35, 40, 45, 50, 55, 60, 65, 70 \\
Feature selector & SelectKBest ($f$-classif) \\
Training adjustments & LOPO-group equalization + SMOTE \\
Validation & Leave-one-participant-out \\
Evaluation level & Patient-level aggregation \\
\bottomrule
\end{tabular}

\end{table}

Baseline recordings were collected during eyes-closed rest with a 4-channel montage (\texttt{AF7}, \texttt{AF8}, \texttt{TP9}, \texttt{TP10}) sampled at 256 Hz. The available study exports comprised 10 s segments with 2,560 samples per channel, and the underlying recording sessions were stored either as a single 10 min file or as two 5 min files depending on participant. Across the cohort, the export set contained 1,203 windows, including 393 remission-labeled and 810 non-remission-labeled windows. This acquisition setup uses a limited frontal-temporal montage with clear portable-deployment relevance \citep{Casson2010,stoyell2021high}. The montage also constrained the downstream feature space toward frontal, temporal, cross-hemispheric, and frontal-temporal interaction measures.

For the present sweep, the processing configuration applied participant-level per-channel demeaning before upsampling, followed by a fourth-order Butterworth low-pass filter at 60 Hz, a downsampling step back to 256 Hz, and zero-overlap window slicing. The sweep then compared non-overlapping analysis windows of 2, 4, 6, 8, and 10 s, so the model inputs included both the native 10 s exports and shorter subdivisions derived from the same recordings. These choices materially affect downstream EEG feature distributions, especially in low-density or portable pipelines \citep{Siuly2017,Miranda2019,Mahato2020,babiloni2012_guidelines,simmatis2023_technical}. The feature extractor enabled spectral, temporal, entropy, complexity, connectivity, coherence, asymmetry, and cross-hemispheric descriptors, including frontal/temporal asymmetry, cross-regional ratios, and inter-channel coupling measures. This broad handcrafted feature pool was intended to preserve interpretability while allowing sparse downstream selection.

\subsection{Sweep design and validation}

The experiment runner launched two model families over the same feature-selection sweep: an advanced hybrid 1D CNN-LSTM and a linear SVM. All runs used sequential window ordering, SelectKBest with ANOVA $F$ statistics (\texttt{select\_k\_best\_f\_classif}), a fixed consensus feature budget of \texttt{outer-k=10}, leave-one-participant-out outer evaluation, leave-one-participant-out training-group equalization, and SMOTE within training folds. The sweep varied the per-fold feature-selection budget (\texttt{inner-k}) across 15 values from 1 to 70 and repeated the search across 5 window sizes, yielding 150 successful runs in total (75 per model).

This design matters for interpretation. The outer evaluation remained subject independent throughout, so no participant contributed both training and test windows to the same outer fold. The \texttt{inner-k} parameter controlled how many features were selected inside each outer training fold, whereas \texttt{outer-k} controlled only the final consensus feature budget after cross-validation. The feature-selection stability panels in the manuscript describe shared pipeline behavior because both classifiers consumed the same selected features for matched sweep settings.

\subsection{Models}

The SVM baseline used a linear kernel with class-balanced weighting and probability output enabled. The hybrid model used a substantially richer architecture defined in the project model configuration: stacked convolutional blocks, bidirectional LSTM layers, multi-head attention with positional encoding, feature-pyramid fusion, and GELU-activated dense layers with dropout and batch normalization. The manuscript treats this architecture as a test case for whether a more expressive learner can exploit a very small selected feature set more effectively than a linear baseline.

\subsection{Outcome metrics and biomarker summaries}

The primary endpoint was patient-level ROC-AUC, computed from participant probabilities aggregated over held-out windows. Secondary patient-level metrics included PR-AUC, balanced accuracy, accuracy, F1, MCC, precision, recall, and specificity. The pipeline also reported bootstrap confidence intervals and a permutation test for the patient-level ROC-AUC summary. All main manuscript claims use patient-level metrics. Window-level scores appear only as intermediate outputs.

To characterize biomarker behavior across outer folds, we used several post hoc summaries: mean pairwise Jaccard overlap, Kuncheva index, unique feature count, top-feature share, class-conditional selection deltas, and effect-direction consistency. These summaries are descriptive. They characterize recurrence and directionality across folds and do not support causal interpretation. Low overlap can coexist with recurring features that retain consistent class-separation direction when selected \citep{olbrich2013_eeg,olbrich2015_personalized,simmatis2023_technical}.
